\documentclass[12pt,a4paper]{article}

% ==================== PACKAGES ====================
\usepackage[utf8]{inputenc}
\usepackage[vietnamese]{babel}
\usepackage{amsmath,amssymb,amsthm}
\usepackage{geometry}
\usepackage{enumitem}
\usepackage[most]{tcolorbox}
\usepackage{hyperref}
\usepackage{fancyhdr}
\usepackage{graphicx}
\usepackage{tikz}
\usepackage{eso-pic}
\usepackage{xcolor}

\geometry{a4paper, margin=2cm, bottom=2.5cm}
\hypersetup{colorlinks=true, linkcolor=blue!70!black, urlcolor=blue}
\setlist[itemize]{topsep=4pt, itemsep=2pt}
\setlist[enumerate]{topsep=4pt, itemsep=2pt}

% ==================== WATERMARK ====================
\AddToShipoutPictureFG{%
  \AtPageCenter{%
    \begin{tikzpicture}[remember picture, overlay]
      \node[rotate=45, scale=1, opacity=0.06]{%
        \fontsize{60}{72}\selectfont\bfseries\sffamily Đặng Tấn Quí%
      };
    \end{tikzpicture}%
  }%
}

% ==================== HEADER / FOOTER ====================
\pagestyle{fancy}
\fancyhf{}
\fancyhead[L]{\small\textit{Đạo hàm -- Định nghĩa và Tính chất}}
\fancyhead[R]{\small\textit{Đặng Tấn Quí}}
\fancyfoot[C]{\thepage}
\fancyfoot[L]{\footnotesize\textcopyright\ Đặng Tấn Quí}
\fancyfoot[R]{\footnotesize SĐT: 0377\,122\,210}
\renewcommand{\headrulewidth}{0.4pt}
\renewcommand{\footrulewidth}{0.4pt}

\fancypagestyle{plain}{%
  \fancyhf{}
  \fancyfoot[C]{\thepage}
  \fancyfoot[L]{\footnotesize\textcopyright\ Đặng Tấn Quí}
  \fancyfoot[R]{\footnotesize SĐT: 0377\,122\,210}
  \renewcommand{\headrulewidth}{0pt}
  \renewcommand{\footrulewidth}{0.4pt}
}

% ==================== CUSTOM ENVIRONMENTS ====================
\newtcolorbox{dinhnghia}[1][]{%
  enhanced, breakable,
  colback=blue!3!white, colframe=blue!60!black,
  fonttitle=\bfseries, title={Định nghĩa}, #1%
}

\newtcolorbox{dinhly}[1][]{%
  enhanced, breakable,
  colback=green!3!white, colframe=green!50!black,
  fonttitle=\bfseries, title={Định lý}, #1%
}

\newtcolorbox{tinhchat}[1][]{%
  enhanced, breakable,
  colback=yellow!5!white, colframe=orange!50!black,
  fonttitle=\bfseries, title={Tính chất}, #1%
}

\newtcolorbox{phuongphap}[1][]{%
  enhanced, breakable,
  colback=gray!5!white, colframe=gray!50!black,
  fonttitle=\bfseries, title={Phương pháp}, #1%
}

\newtcolorbox{luuy}[1][]{%
  enhanced, breakable,
  colback=red!3!white, colframe=red!50!black,
  fonttitle=\bfseries, title={Lưu ý}, #1%
}

\newtcolorbox{congthuc}[1][]{%
  enhanced, breakable,
  colback=cyan!3!white, colframe=cyan!50!black,
  fonttitle=\bfseries, title={Công thức}, #1%
}

\newtcolorbox{vidu}[1][]{%
  enhanced, breakable,
  colback=violet!3!white, colframe=violet!50!black,
  fonttitle=\bfseries, title={Ví dụ}, #1%
}

% ==================== DOCUMENT ====================
\begin{document}

% ==================== TRANG BÌA ====================
\thispagestyle{empty}
\begin{tikzpicture}[remember picture, overlay]
  \fill[blue!80!black] (current page.south west) rectangle (current page.north east);
  \fill[blue!60!cyan, opacity=0.8]
    (current page.south west) -- (current page.north west) --
    ([xshift=-3cm]current page.north east) -- (current page.south east) -- cycle;

  \foreach \r/\o in {6/0.05, 4.5/0.08, 3/0.12} {
    \fill[white, opacity=\o] ([xshift=2cm, yshift=-2cm]current page.north east) circle (\r cm);
  }

  \draw[white, line width=2pt]
    ([yshift=-5cm]current page.north west) ++(2cm,0) -- ++(17cm,0);
  \draw[white, line width=2pt]
    ([yshift=-16cm]current page.north west) ++(2cm,0) -- ++(17cm,0);

  \node[anchor=west, white, font=\fontsize{30}{36}\bfseries\selectfont]
    at ([xshift=3cm, yshift=-7.5cm]current page.north west)
    {PHẦN 2: ĐẠO HÀM};
  \node[anchor=west, white, font=\fontsize{22}{28}\bfseries\selectfont]
    at ([xshift=3cm, yshift=-9.2cm]current page.north west)
    {ĐỊNH NGHĨA VÀ TÍNH CHẤT};

  \node[anchor=west, white, font=\fontsize{16}{20}\selectfont]
    at ([xshift=3cm, yshift=-11cm]current page.north west)
    {Tài liệu luyện thi du học};

  \node[white, opacity=0.12, font=\fontsize{120}{120}\selectfont, rotate=-15]
    at ([xshift=-3cm, yshift=4cm]current page.center)
    {$f'$};
  \node[white, opacity=0.10, font=\fontsize{80}{80}\selectfont, rotate=10]
    at ([xshift=5cm, yshift=3cm]current page.center)
    {$\dfrac{dy}{dx}$};

  \node[anchor=west, white, font=\Large\bfseries]
    at ([xshift=3cm, yshift=-13cm]current page.north west)
    {Biên soạn:};
  \node[anchor=west, yellow!80!white, font=\fontsize{28}{34}\bfseries\selectfont]
    at ([xshift=3cm, yshift=-14.5cm]current page.north west)
    {ĐẶNG TẤN QUÍ};

  \node[anchor=west, white!90, font=\large]
    at ([xshift=3cm, yshift=-18cm]current page.north west)
    {SĐT: 0377\,122\,210};

  \node[anchor=south, white!80, font=\small]
    at ([yshift=1.5cm]current page.south)
    {Tài liệu có bản quyền --- Cấm sao chép, phân phối khi chưa được sự đồng ý của tác giả};

  \node[anchor=south east, white, font=\Large\bfseries]
    at ([xshift=-2cm, yshift=3cm]current page.south east)
    {\the\year};
\end{tikzpicture}

\newpage

\thispagestyle{empty}
\vspace*{5cm}
\begin{center}
  {\Large\bfseries PHẦN 2: ĐẠO HÀM -- ĐỊNH NGHĨA VÀ TÍNH CHẤT}\\[8pt]
  {\large Tài liệu luyện thi du học}
\end{center}
\vspace{2cm}
\begin{tcolorbox}[enhanced, colback=red!3!white, colframe=red!60!black, fonttitle=\bfseries, title={Thông tin bản quyền}, boxrule=1.5pt]
  \begin{itemize}[leftmargin=*]
    \item \textbf{Tác giả:} Đặng Tấn Quí
    \item \textbf{Liên hệ:} 0377\,122\,210
    \item \textbf{Năm:} \the\year
  \end{itemize}
  \vspace{8pt}
  \textbf{Bản quyền \textcopyright\ \the\year\ Đặng Tấn Quí. Bảo lưu mọi quyền.}

  \vspace{6pt}
  Nghiêm cấm mọi hình thức sao chép, tái bản, phân phối dưới bất kỳ hình thức nào mà không có sự đồng ý bằng văn bản của tác giả.
\end{tcolorbox}
\vfill
\begin{center}
  \small\textit{Tài liệu lưu hành nội bộ --- Không phát hành thương mại}
\end{center}

\newpage
\tableofcontents
\newpage

%% ================================================================
%%  PHẦN 2: ĐẠO HÀM -- ĐỊNH NGHĨA VÀ TÍNH CHẤT
%% ================================================================
\section{Đạo hàm -- Định nghĩa và Tính chất}

%% ----------------------------------------------------------------
%%  2.1 Định nghĩa đạo hàm qua giới hạn
%% ----------------------------------------------------------------
\subsection{Định nghĩa đạo hàm qua giới hạn (Difference Quotient)}

\begin{dinhnghia}[title={Đạo hàm tại một điểm}]
  Cho hàm $y = f(x)$. \textbf{Đạo hàm của $f$ tại $x = a$}, ký hiệu $f'(a)$, được định nghĩa bởi:
  \[
    f'(a) = \lim_{h \to 0} \dfrac{f(a + h) - f(a)}{h},
  \]
  miễn là giới hạn này tồn tại và hữu hạn. Biểu thức $\dfrac{f(a+h) - f(a)}{h}$ gọi là \textbf{thương sai phân (difference quotient)}.

  \medskip
  \textbf{Dạng tương đương:}
  \[
    f'(a) = \lim_{x \to a} \dfrac{f(x) - f(a)}{x - a}.
  \]
\end{dinhnghia}

\begin{dinhnghia}[title={Hàm đạo hàm}]
  \textbf{Hàm đạo hàm} $f'(x)$ (hay $\dfrac{dy}{dx}$) là hàm số xác định bởi:
  \[
    f'(x) = \lim_{h \to 0} \dfrac{f(x + h) - f(x)}{h},
  \]
  tại mọi điểm $x$ mà giới hạn trên tồn tại.

  \medskip
  \textbf{Các ký hiệu thông dụng:} $f'(x)$, $y'$, $\dfrac{dy}{dx}$, $\dfrac{df}{dx}$, $Df(x)$.
\end{dinhnghia}

\begin{dinhnghia}[title={Ý nghĩa hình học}]
  $f'(a)$ là \textbf{hệ số góc (slope) của tiếp tuyến} với đồ thị $y = f(x)$ tại điểm $M\bigl(a,\, f(a)\bigr)$.
  \[
    \text{Phương trình tiếp tuyến tại } M: \quad y = f'(a)(x - a) + f(a).
  \]
\end{dinhnghia}

\begin{dinhnghia}[title={Ý nghĩa vật lý}]
  Nếu $s(t)$ là hàm vị trí của một vật theo thời gian, thì $s'(t)$ là \textbf{vận tốc tức thời} của vật tại thời điểm $t$.
\end{dinhnghia}

\begin{vidu}
  Tìm $f'(x)$ từ định nghĩa với $f(x) = x^2$.

  \medskip\textbf{Giải:}
  \[
    f'(x) = \lim_{h \to 0} \dfrac{(x+h)^2 - x^2}{h}
    = \lim_{h \to 0} \dfrac{x^2 + 2xh + h^2 - x^2}{h}
    = \lim_{h \to 0} \dfrac{2xh + h^2}{h}
    = \lim_{h \to 0} (2x + h) = 2x.
  \]
\end{vidu}

\medskip\noindent\textbf{Các dạng bài tập:}
\begin{enumerate}[label=\textbf{Dạng \arabic*:}, leftmargin=*, align=left]
  \item Tính đạo hàm từ định nghĩa giới hạn.
  \item Tính hệ số góc tiếp tuyến và viết phương trình tiếp tuyến.
  \item Tìm vận tốc tức thời từ hàm vị trí.
\end{enumerate}

%% ----------------------------------------------------------------
%%  2.2 Tính khả vi và tính liên tục
%% ----------------------------------------------------------------
\subsection{Mối liên hệ giữa tính khả vi (Differentiability) và tính liên tục}

\begin{dinhly}[title={Khả vi kéo theo liên tục}]
  Nếu $f$ \textbf{khả vi} (có đạo hàm) tại $x = a$, thì $f$ \textbf{liên tục} tại $x = a$.

  \medskip
  \textbf{Ngược lại không đúng:} Hàm liên tục không nhất thiết khả vi.
\end{dinhly}

\begin{dinhnghia}[title={Các trường hợp hàm không khả vi}]
  Hàm $f$ \textbf{không khả vi} tại $x = a$ trong các trường hợp:
  \begin{enumerate}
    \item \textbf{Điểm góc (Corner):} Đồ thị có ``góc nhọn''; giới hạn trái và phải của thương sai phân tồn tại nhưng khác nhau. Ví dụ: $f(x) = |x|$ tại $x = 0$.
    \item \textbf{Tiếp tuyến đứng (Vertical tangent):} $f'(a) = \pm\infty$. Ví dụ: $f(x) = \sqrt[3]{x}$ tại $x = 0$.
    \item \textbf{Điểm gián đoạn:} Hàm không liên tục tại $a$ $\Rightarrow$ không khả vi tại $a$.
    \item \textbf{Điểm nhọn (Cusp):} Tương tự góc, hai đường tiếp tuyến có hệ số góc khác nhau.
  \end{enumerate}
\end{dinhnghia}

\begin{vidu}
  Chứng minh $f(x) = |x|$ không khả vi tại $x = 0$, dù liên tục tại đó.

  \medskip\textbf{Giải:}
  \begin{itemize}
    \item $\displaystyle\lim_{h \to 0^+} \dfrac{|0+h| - |0|}{h} = \lim_{h \to 0^+} \dfrac{h}{h} = 1$.
    \item $\displaystyle\lim_{h \to 0^-} \dfrac{|0+h| - |0|}{h} = \lim_{h \to 0^-} \dfrac{-h}{h} = -1$.
  \end{itemize}
  Vì giới hạn trái $\ne$ giới hạn phải, $f'(0)$ không tồn tại. Tuy nhiên $f(x) = |x|$ rõ ràng liên tục tại $0$.
\end{vidu}

%% ----------------------------------------------------------------
%%  2.3 Các quy tắc đạo hàm cơ bản
%% ----------------------------------------------------------------
\subsection{Các quy tắc đạo hàm cơ bản: Power Rule, Constant, Sum/Difference}

\begin{congthuc}[title={Bảng đạo hàm cơ bản}]
  \begin{align*}
    &\dfrac{d}{dx}[c] = 0 \quad (c \text{ là hằng số})
    &\qquad&
    \dfrac{d}{dx}[x^n] = n x^{n-1} \quad (\text{Power Rule}) \\[6pt]
    &\dfrac{d}{dx}[x] = 1
    &&
    \dfrac{d}{dx}[\sqrt{x}] = \dfrac{1}{2\sqrt{x}} \\[6pt]
    &\dfrac{d}{dx}\!\left[\dfrac{1}{x}\right] = -\dfrac{1}{x^2}
    &&
    \dfrac{d}{dx}[x^{1/n}] = \dfrac{1}{n} x^{1/n - 1}
  \end{align*}
\end{congthuc}

\begin{tinhchat}[title={Các quy tắc kết hợp}]
  \begin{itemize}
    \item \textbf{Nhân hằng số:} $\dfrac{d}{dx}[c \cdot f(x)] = c \cdot f'(x)$.
    \item \textbf{Tổng/Hiệu:} $\dfrac{d}{dx}[f(x) \pm g(x)] = f'(x) \pm g'(x)$.
  \end{itemize}
\end{tinhchat}

\begin{vidu}
  Tính đạo hàm của $f(x) = 3x^4 - 5x^2 + 7x - 2$.

  \medskip\textbf{Giải:}
  \[
    f'(x) = 3 \cdot 4x^3 - 5 \cdot 2x + 7 = 12x^3 - 10x + 7.
  \]
\end{vidu}

%% ----------------------------------------------------------------
%%  2.4 Đạo hàm các hàm lượng giác
%% ----------------------------------------------------------------
\subsection{Đạo hàm các hàm số lượng giác}

\begin{congthuc}[title={Đạo hàm hàm lượng giác}]
  \begin{align*}
    &\dfrac{d}{dx}[\sin x] = \cos x
    &\qquad&
    \dfrac{d}{dx}[\cos x] = -\sin x \\[6pt]
    &\dfrac{d}{dx}[\tan x] = \dfrac{1}{\cos^2 x}
    &&
    \dfrac{d}{dx}[\cot x] = -\dfrac{1}{\sin^2 x}
  \end{align*}
\end{congthuc}

\begin{luuy}
  Các công thức trên được chứng minh từ định nghĩa giới hạn, dùng các giới hạn đặc biệt $\displaystyle\lim_{h \to 0} \dfrac{\sin h}{h} = 1$ và $\displaystyle\lim_{h \to 0} \dfrac{\cos h - 1}{h} = 0$.
\end{luuy}

\begin{vidu}
  Tính đạo hàm của $f(x) = 3\sin x - 2\cos x + \tan x$.

  \medskip\textbf{Giải:}
  \[
    f'(x) = 3\cos x + 2\sin x + \sec^2 x.
  \]
\end{vidu}

%% ----------------------------------------------------------------
%%  2.5 Quy tắc Tích và Quy tắc Thương
%% ----------------------------------------------------------------
\subsection{Quy tắc Tích (Product Rule) và Quy tắc Thương (Quotient Rule)}

\begin{congthuc}[title={Quy tắc Tích (Product Rule)}]
  Nếu $f$ và $g$ khả vi thì:
  \[
    \dfrac{d}{dx}[f(x) \cdot g(x)] = f'(x) \cdot g(x) + f(x) \cdot g'(x).
  \]
\end{congthuc}

\begin{congthuc}[title={Quy tắc Thương (Quotient Rule)}]
  Nếu $f$ và $g$ khả vi và $g(x) \ne 0$ thì:
  \[
    \dfrac{d}{dx}\!\left[\dfrac{f(x)}{g(x)}\right] = \dfrac{f'(x) \cdot g(x) - f(x) \cdot g'(x)}{[g(x)]^2}.
  \]
\end{congthuc}

\begin{vidu}
  \begin{enumerate}[label=(\alph*)]
    \item Tính đạo hàm của $h(x) = x^2 \sin x$.
    \item Tính đạo hàm của $k(x) = \dfrac{x^2 + 1}{x - 3}$.
  \end{enumerate}

  \medskip\textbf{Giải:}
  \begin{enumerate}[label=(\alph*)]
    \item $h'(x) = 2x \sin x + x^2 \cos x$.
    \item $k'(x) = \dfrac{2x(x-3) - (x^2 + 1)(1)}{(x-3)^2} = \dfrac{2x^2 - 6x - x^2 - 1}{(x-3)^2} = \dfrac{x^2 - 6x - 1}{(x-3)^2}$.
  \end{enumerate}
\end{vidu}

\medskip\noindent\textbf{Các dạng bài tập:}
\begin{enumerate}[label=\textbf{Dạng \arabic*:}, leftmargin=*, align=left]
  \item Tính đạo hàm của tích hai hàm số.
  \item Tính đạo hàm của thương hai hàm số.
  \item Kết hợp Product Rule và Quotient Rule.
\end{enumerate}

%% ----------------------------------------------------------------
%%  2.6 Quy tắc Hàm hợp (Chain Rule)
%% ----------------------------------------------------------------
\subsection{Quy tắc Hàm hợp (Chain Rule) -- Trọng tâm quan trọng nhất}

\begin{dinhly}[title={Chain Rule}]
  Nếu $g$ khả vi tại $x$ và $f$ khả vi tại $g(x)$, thì hàm hợp $h(x) = f(g(x))$ khả vi tại $x$ và:
  \[
    h'(x) = f'\bigl(g(x)\bigr) \cdot g'(x).
  \]
  \textbf{Ký hiệu Leibniz:} Nếu $y = f(u)$ và $u = g(x)$, thì:
  \[
    \dfrac{dy}{dx} = \dfrac{dy}{du} \cdot \dfrac{du}{dx}.
  \]
\end{dinhly}

\begin{phuongphap}[title={Nhận dạng và áp dụng Chain Rule}]
  \begin{enumerate}
    \item Xác định hàm ``ngoài'' $f$ và hàm ``trong'' $u = g(x)$.
    \item Tính $f'(u)$ (đạo hàm của ngoài theo $u$).
    \item Tính $g'(x)$ (đạo hàm của trong theo $x$).
    \item Kết quả: $h'(x) = f'(g(x)) \cdot g'(x)$.
  \end{enumerate}

  \textbf{Bảng nhận dạng nhanh:}
  \begin{itemize}
    \item $[\sin(g(x))]' = \cos(g(x)) \cdot g'(x)$
    \item $[\cos(g(x))]' = -\sin(g(x)) \cdot g'(x)$
    \item $[g(x)^n]' = n[g(x)]^{n-1} \cdot g'(x)$
    \item $[e^{g(x)}]' = e^{g(x)} \cdot g'(x)$
    \item $[\ln(g(x))]' = \dfrac{g'(x)}{g(x)}$
    \item $[\sqrt{g(x)}]' = \dfrac{g'(x)}{2\sqrt{g(x)}}$
  \end{itemize}
\end{phuongphap}

\begin{vidu}
  Tính đạo hàm của:
  \begin{enumerate}[label=(\alph*)]
    \item $f(x) = \sin(3x^2 + 1)$
    \item $g(x) = (x^3 - 5x)^{10}$
    \item $h(x) = \sqrt{2x^2 - 1}$
    \item $k(x) = \cos^3(x) = [\cos x]^3$
  \end{enumerate}

  \medskip\textbf{Giải:}
  \begin{enumerate}[label=(\alph*)]
    \item $f'(x) = \cos(3x^2 + 1) \cdot 6x$.
    \item $g'(x) = 10(x^3 - 5x)^9 \cdot (3x^2 - 5)$.
    \item $h'(x) = \dfrac{4x}{2\sqrt{2x^2 - 1}} = \dfrac{2x}{\sqrt{2x^2 - 1}}$.
    \item $k'(x) = 3\cos^2 x \cdot (-\sin x) = -3\cos^2 x \sin x$.
  \end{enumerate}
\end{vidu}

\begin{vidu}[title={Chain Rule nhiều lớp}]
  Tính đạo hàm của $f(x) = \sin^2(3x)$.

  \medskip\textbf{Giải:}
  Nhận dạng: $f(x) = [\sin(3x)]^2$ -- có ba lớp.
  \[
    f'(x) = 2\sin(3x) \cdot \cos(3x) \cdot 3 = 6\sin(3x)\cos(3x) = 3\sin(6x).
  \]
\end{vidu}

\medskip\noindent\textbf{Các dạng bài tập:}
\begin{enumerate}[label=\textbf{Dạng \arabic*:}, leftmargin=*, align=left]
  \item Chain Rule với hàm lũy thừa, lượng giác, căn thức.
  \item Chain Rule nhiều lớp (hàm hợp của hàm hợp).
  \item Kết hợp Chain Rule với Product Rule hoặc Quotient Rule.
  \item Tìm phương trình tiếp tuyến khi dùng Chain Rule.
\end{enumerate}

%% ----------------------------------------------------------------
%%  2.7 Đạo hàm ẩn (Implicit Differentiation)
%% ----------------------------------------------------------------
\subsection{Đạo hàm ẩn (Implicit Differentiation)}

\begin{phuongphap}[title={Phương pháp đạo hàm ẩn}]
  Dùng khi $y$ được xác định \textbf{ẩn} (không giải được $y$ theo $x$ một cách tường minh).
  \begin{enumerate}
    \item Đạo hàm cả hai vế theo $x$, coi $y$ là hàm của $x$.
    \item Mỗi khi đạo hàm $y$, áp dụng Chain Rule: $\dfrac{d}{dx}[f(y)] = f'(y) \cdot \dfrac{dy}{dx}$.
    \item Tập hợp tất cả các số hạng chứa $\dfrac{dy}{dx}$ về một phía.
    \item Giải phương trình để tìm $\dfrac{dy}{dx}$.
  \end{enumerate}
\end{phuongphap}

\begin{vidu}
  Tìm $\dfrac{dy}{dx}$ với $x^2 + y^2 = 25$ (đường tròn tâm $O$, bán kính $5$).

  \medskip\textbf{Giải:}
  Đạo hàm hai vế theo $x$:
  \[
    2x + 2y \cdot \dfrac{dy}{dx} = 0
    \;\Rightarrow\;
    \dfrac{dy}{dx} = -\dfrac{x}{y} \quad (y \ne 0).
  \]
\end{vidu}

\begin{vidu}
  Tìm phương trình tiếp tuyến với $x^3 + y^3 = 6xy$ tại điểm $(3, 3)$.

  \medskip\textbf{Giải:}
  Đạo hàm ẩn:
  \[
    3x^2 + 3y^2 y' = 6y + 6x y'
    \;\Rightarrow\;
    y'(3y^2 - 6x) = 6y - 3x^2
    \;\Rightarrow\;
    y' = \dfrac{6y - 3x^2}{3y^2 - 6x}.
  \]
  Tại $(3, 3)$: $y' = \dfrac{18 - 27}{27 - 18} = \dfrac{-9}{9} = -1$.

  Phương trình tiếp tuyến: $y - 3 = -1(x - 3) \Rightarrow y = -x + 6$.
\end{vidu}

\medskip\noindent\textbf{Các dạng bài tập:}
\begin{enumerate}[label=\textbf{Dạng \arabic*:}, leftmargin=*, align=left]
  \item Tìm $\dfrac{dy}{dx}$ bằng đạo hàm ẩn.
  \item Tìm phương trình tiếp tuyến của đường cong ẩn.
  \item Tính đạo hàm bậc hai $\dfrac{d^2y}{dx^2}$ bằng đạo hàm ẩn.
\end{enumerate}

%% ----------------------------------------------------------------
%%  2.8 Đạo hàm hàm mũ và logarit
%% ----------------------------------------------------------------
\subsection{Đạo hàm của hàm mũ và hàm Logarit}

\begin{congthuc}[title={Đạo hàm hàm mũ}]
  \begin{align*}
    &\dfrac{d}{dx}[e^x] = e^x
    &\qquad&
    \dfrac{d}{dx}[a^x] = a^x \ln a \quad (a > 0,\, a \ne 1) \\[6pt]
    &\dfrac{d}{dx}[e^{u(x)}] = e^{u(x)} \cdot u'(x)
    &&
    \dfrac{d}{dx}[a^{u(x)}] = a^{u(x)} \ln a \cdot u'(x)
  \end{align*}
\end{congthuc}

\begin{congthuc}[title={Đạo hàm hàm Logarit}]
  \begin{align*}
    &\dfrac{d}{dx}[\ln x] = \dfrac{1}{x} \quad (x > 0)
    &\qquad&
    \dfrac{d}{dx}[\log_a x] = \dfrac{1}{x \ln a} \quad (x > 0) \\[6pt]
    &\dfrac{d}{dx}[\ln|x|] = \dfrac{1}{x}
    &&
    \dfrac{d}{dx}[\ln u(x)] = \dfrac{u'(x)}{u(x)}
  \end{align*}
\end{congthuc}

\begin{phuongphap}[title={Logarithmic Differentiation (Đạo hàm Logarit)}]
  Dùng khi hàm số có dạng $y = [f(x)]^{g(x)}$ (cơ số và số mũ đều phụ thuộc $x$):
  \begin{enumerate}
    \item Lấy $\ln$ cả hai vế: $\ln y = g(x) \ln f(x)$.
    \item Đạo hàm ẩn hai vế: $\dfrac{y'}{y} = g'(x) \ln f(x) + g(x) \cdot \dfrac{f'(x)}{f(x)}$.
    \item Nhân hai vế với $y$.
  \end{enumerate}
\end{phuongphap}

\begin{vidu}
  Tính đạo hàm của $y = x^x$ ($x > 0$).

  \medskip\textbf{Giải:}
  $\ln y = x \ln x$. Đạo hàm hai vế:
  \[
    \dfrac{y'}{y} = \ln x + x \cdot \dfrac{1}{x} = \ln x + 1.
    \;\Rightarrow\;
    y' = y(\ln x + 1) = x^x(\ln x + 1).
  \]
\end{vidu}

\medskip\noindent\textbf{Các dạng bài tập:}
\begin{enumerate}[label=\textbf{Dạng \arabic*:}, leftmargin=*, align=left]
  \item Tính đạo hàm hàm mũ cơ số $e$ và $a$.
  \item Tính đạo hàm hàm logarit.
  \item Kết hợp Chain Rule với hàm mũ và logarit.
  \item Đạo hàm logarit để tính $[f(x)]^{g(x)}$.
\end{enumerate}

%% ----------------------------------------------------------------
%%  2.9 Đạo hàm hàm ngược và hàm lượng giác ngược
%% ----------------------------------------------------------------
\subsection{Đạo hàm của hàm ngược và hàm lượng giác ngược}

\begin{dinhly}[title={Đạo hàm hàm ngược}]
  Nếu $f$ khả vi và có hàm ngược $f^{-1}$, thì tại điểm $y = f(x)$:
  \[
    (f^{-1})'(y) = \dfrac{1}{f'(x)} = \dfrac{1}{f'(f^{-1}(y))}.
  \]
  \textbf{Ý nghĩa:} Hệ số góc của tiếp tuyến tại $(y, x)$ trên đồ thị $f^{-1}$ bằng nghịch đảo hệ số góc tại $(x, y)$ trên đồ thị $f$.
\end{dinhly}

\begin{congthuc}[title={Đạo hàm hàm lượng giác ngược}]
  \begin{align*}
    &\dfrac{d}{dx}[\arcsin x] = \dfrac{1}{\sqrt{1 - x^2}}, \quad |x| < 1
    &\qquad&
    \dfrac{d}{dx}[\arccos x] = -\dfrac{1}{\sqrt{1 - x^2}}, \quad |x| < 1 \\[6pt]
    &\dfrac{d}{dx}[\arctan x] = \dfrac{1}{1 + x^2}
    &&
    \dfrac{d}{dx}[\text{arccot}\, x] = -\dfrac{1}{1 + x^2}
  \end{align*}
\end{congthuc}

\begin{luuy}
  \begin{itemize}
    \item Kết hợp với Chain Rule: $\dfrac{d}{dx}[\arcsin(u)] = \dfrac{u'}{\sqrt{1-u^2}}$, $\dfrac{d}{dx}[\arctan(u)] = \dfrac{u'}{1+u^2}$.
    \item Nhớ: $\arcsin$ và $\arccos$ bổ sung nhau: đạo hàm của chúng đối nhau (một cái dương, một cái âm).
  \end{itemize}
\end{luuy}

\begin{vidu}
  Tính đạo hàm của:
  \begin{enumerate}[label=(\alph*)]
    \item $f(x) = \arctan(2x)$
    \item $g(x) = \arcsin(x^2)$
    \item $h(x) = x \arctan x - \dfrac{1}{2}\ln(1 + x^2)$
  \end{enumerate}

  \medskip\textbf{Giải:}
  \begin{enumerate}[label=(\alph*)]
    \item $f'(x) = \dfrac{2}{1 + (2x)^2} = \dfrac{2}{1 + 4x^2}$.
    \item $g'(x) = \dfrac{2x}{\sqrt{1 - x^4}}$.
    \item $h'(x) = \arctan x + \dfrac{x}{1+x^2} - \dfrac{x}{1+x^2} = \arctan x$.
  \end{enumerate}
\end{vidu}

\medskip\noindent\textbf{Các dạng bài tập:}
\begin{enumerate}[label=\textbf{Dạng \arabic*:}, leftmargin=*, align=left]
  \item Tính đạo hàm hàm lượng giác ngược cơ bản.
  \item Kết hợp Chain Rule với lượng giác ngược.
  \item Tính đạo hàm hàm hợp phức tạp chứa arcsin, arccos, arctan.
  \item Bài toán tìm tiếp tuyến, pháp tuyến với đồ thị hàm lượng giác ngược.
\end{enumerate}

\end{document}
